% ====================================================================
% FF2-3 (ch1 §9 fig:reversal 다음 배치 제안) — L_V 성장에 따른 peak 변형·이웃 peak 겹침
%   좌표 = 인과 기억 적분(eq:lag)·peak shape(eq:peakshape)의 점별 수치 평가(python; 방전,
%   w=1 단위): P(V;L)=∫_0^∞ e^{-t}·ξ_eq(1−ξ_eq)|_{V−Lt} dt.
%   검산: L=1.5w 정점 (+1.01w, 0.1955) = fig:reversal 과 일치. 면적은 모든 L 에서 1(=Q_j 보존).
%   (b) 두 전이(간격 6w, 같은 Q)의 합: L=0.3w 골/정점=0.38 → L=3w 0.83 (겹침).
%   사용 라이브러리: ch1 preamble 범위 내(arrows.meta).
% ====================================================================
\begin{figure}[t]
\centering
\begin{tikzpicture}
% =================== (a) single transition, L family ===================
\begin{scope}[x=0.40cm,y=13.5cm]
\draw[-{Latex[length=1.6mm]}] (-6.3,0) -- (-6.3,0.285) node[above,font=\scriptsize] {peak [$1/w$]};
\draw[-{Latex[length=1.6mm]}] (-6.3,0) -- (10.6,0) node[below,font=\scriptsize] {$(V-U_j)/w$};
\foreach \xx in {-4,0,4,8} {\draw (\xx,0.004) -- (\xx,-0.004) node[below,font=\scriptsize] {\xx};}
\foreach \yy/\lab in {0.1/0.1,0.2/0.2,0.25/0.25}
  {\draw (-6.15,\yy) -- (-6.45,\yy) node[left,font=\scriptsize] {\lab};}
% L=0 equilibrium bell (dotted)
\draw[densely dotted,thick] plot[smooth] coordinates {(-6.00,0.0025) (-5.42,0.0044) (-4.84,0.0078) (-4.25,0.0138) (-3.67,0.0242) (-3.09,0.0416) (-2.51,0.0696) (-1.93,0.1109) (-1.35,0.1639) (-0.76,0.2168) (-0.18,0.2479) (0.00,0.2500) (0.40,0.2403) (0.98,0.1983) (1.56,0.1432) (2.15,0.0938) (2.73,0.0576) (3.31,0.0340) (3.89,0.0196) (4.47,0.0112) (5.05,0.0063) (5.64,0.0035) (6.22,0.0020) (7.38,0.0006) (8.55,0.0002) (9.71,0.0001)};
% L=0.75 (gray solid)
\draw[black!55,thick] plot[smooth] coordinates {(-6.00,0.0014) (-5.13,0.0034) (-4.25,0.0080) (-3.38,0.0185) (-2.95,0.0280) (-2.51,0.0416) (-2.07,0.0609) (-1.64,0.0865) (-1.20,0.1185) (-0.76,0.1544) (-0.33,0.1890) (0.11,0.2153) (0.55,0.2269) (0.98,0.2215) (1.42,0.2013) (1.85,0.1719) (2.29,0.1393) (2.73,0.1081) (3.16,0.0811) (3.60,0.0592) (4.04,0.0422) (4.47,0.0296) (4.91,0.0205) (5.35,0.0140) (5.78,0.0095) (6.22,0.0064) (7.09,0.0028) (7.96,0.0012) (8.84,0.0005) (9.71,0.0002)};
% L=1.5 (solid)
\draw[thick] plot[smooth] coordinates {(-6.00,0.0010) (-5.13,0.0024) (-4.25,0.0056) (-3.38,0.0130) (-2.95,0.0197) (-2.51,0.0295) (-2.07,0.0434) (-1.64,0.0623) (-1.20,0.0864) (-0.76,0.1148) (-0.33,0.1445) (0.11,0.1707) (0.55,0.1887) (0.98,0.1955) (1.42,0.1908) (1.85,0.1768) (2.29,0.1571) (2.73,0.1350) (3.16,0.1128) (3.60,0.0923) (4.04,0.0743) (4.47,0.0590) (4.91,0.0464) (5.35,0.0362) (5.78,0.0280) (6.22,0.0216) (7.09,0.0126) (7.96,0.0073) (8.84,0.0042) (9.71,0.0024)};
% L=3.0 (thick dashed)
\draw[thick,densely dashed] plot[smooth] coordinates {(-6.00,0.0006) (-5.13,0.0015) (-4.25,0.0035) (-3.38,0.0082) (-2.95,0.0124) (-2.51,0.0186) (-2.07,0.0275) (-1.64,0.0397) (-1.20,0.0556) (-0.76,0.0750) (-0.33,0.0962) (0.11,0.1168) (0.55,0.1338) (0.98,0.1449) (1.42,0.1493) (1.85,0.1475) (2.29,0.1410) (2.73,0.1313) (3.16,0.1200) (3.60,0.1080) (4.04,0.0963) (4.47,0.0851) (4.91,0.0748) (5.35,0.0655) (5.78,0.0571) (6.22,0.0497) (7.09,0.0375) (7.96,0.0282) (8.84,0.0211) (9.71,0.0158)};
% maxima trajectory
\fill (0,0.2500) circle (1.2pt); \fill (0.62,0.2272) circle (1.2pt);
\fill (1.01,0.1955) circle (1.2pt); \fill (1.50,0.1494) circle (1.2pt);
\draw[-{Latex[length=1.4mm]},black!60] plot[smooth] coordinates {(0.05,0.2495) (0.62,0.2272) (1.01,0.1955) (1.45,0.1520)};
\node[font=\scriptsize,anchor=west,align=left] at (2.30,0.238) {$L_V\!\uparrow$ (저온$\cdot$고율):\\정점$\downarrow\cdot$중심 이동$\cdot$꼬리 연장};
% legend
\node[font=\tiny,anchor=west,align=left] at (-6.0,0.265)
  {점선 $L_V{=}0$(평형 종, 정점 $0.250$@$0$)\ \ 회색 $0.75w$($0.227$@$+0.62w$)\\
   실선 $1.5w$($0.196$@$+1.01w$)\ \ 파선 $3w$($0.149$@$+1.50w$)};
\node[font=\scriptsize] at (2.0,-0.038) {(a)};
\end{scope}
% =================== (b) two transitions, overlap ===================
\begin{scope}[xshift=8.05cm,x=0.34cm,y=13.5cm]
\draw[-{Latex[length=1.6mm]}] (-5.3,0) -- (-5.3,0.285) node[above,font=\scriptsize] {$\sum$ peak [$1/w$]};
\draw[-{Latex[length=1.6mm]}] (-5.3,0) -- (14.4,0) node[below,font=\scriptsize] {$(V-U_A)/w$};
\foreach \xx in {-4,0,4,8,12} {\draw (\xx,0.004) -- (\xx,-0.004) node[below,font=\scriptsize] {\xx};}
\foreach \yy/\lab in {0.1/0.1,0.2/0.2} {\draw (-5.15,\yy) -- (-5.45,\yy) node[left,font=\scriptsize] {\lab};}
% centers
\draw[densely dotted,black!55] (0,0) -- (0,0.262);
\draw[densely dotted,black!55] (6,0) -- (6,0.262);
\node[font=\tiny,black!55,anchor=south] at (0,0.262) {$U_A$};
\node[font=\tiny,black!55,anchor=south] at (6,0.262) {$U_B{=}U_A{+}6w$};
% sum at L=0.3 (solid)
\draw[thick] plot[smooth] coordinates {(-5.00,0.0051) (-4.12,0.0122) (-3.25,0.0282) (-2.81,0.0423) (-2.37,0.0624) (-1.93,0.0897) (-1.49,0.1244) (-1.05,0.1644) (-0.62,0.2040) (-0.18,0.2345) (0.26,0.2474) (0.70,0.2391) (1.14,0.2133) (1.58,0.1784) (2.02,0.1437) (2.45,0.1158) (2.89,0.0985) (3.33,0.0933) (3.77,0.1005) (4.21,0.1198) (4.65,0.1495) (5.08,0.1853) (5.52,0.2195) (5.96,0.2426) (6.40,0.2466) (6.84,0.2297) (7.28,0.1970) (7.72,0.1573) (8.15,0.1184) (8.59,0.0852) (9.03,0.0593) (9.47,0.0403) (9.91,0.0269) (10.78,0.0116) (11.66,0.0049) (12.54,0.0021) (13.42,0.0009)};
% sum at L=3.0 (thick dashed)
\draw[thick,densely dashed] plot[smooth] coordinates {(-5.00,0.0017) (-4.12,0.0040) (-3.25,0.0093) (-2.81,0.0141) (-2.37,0.0212) (-1.93,0.0311) (-1.49,0.0447) (-1.05,0.0620) (-0.62,0.0824) (-0.18,0.1041) (0.26,0.1241) (0.70,0.1397) (1.14,0.1492) (1.58,0.1523) (2.02,0.1501) (2.45,0.1446) (2.89,0.1378) (3.33,0.1315) (3.77,0.1273) (4.21,0.1267) (4.65,0.1305) (5.08,0.1389) (5.52,0.1508) (5.96,0.1641) (6.40,0.1756) (6.84,0.1827) (7.28,0.1839) (7.72,0.1792) (8.15,0.1699) (8.59,0.1574) (9.03,0.1433) (9.47,0.1287) (9.91,0.1145) (10.78,0.0887) (11.66,0.0676) (12.54,0.0510) (13.42,0.0383)};
% valley annotations
\fill (3.33,0.0933) circle (1.1pt);
\draw[-{Latex[length=1.2mm]},black!60] (3.1,0.036) -- (3.30,0.088);
\node[font=\tiny,anchor=north,align=center] at (3.15,0.034) {골$/$정점\\$0.38$ (분리)};
\fill (4.21,0.1267) circle (1.1pt);
\draw[-{Latex[length=1.2mm]},black!60] (5.6,0.075) -- (4.32,0.123);
\node[font=\tiny,anchor=west,align=left] at (5.35,0.062) {골$/$정점 $0.83$\\(겹침 --- 골 메움)};
% legend
\draw[thick] (7.9,0.252) -- (8.9,0.252);
\node[font=\tiny,anchor=west] at (9.0,0.252) {$L_V{=}0.3w$};
\draw[thick,densely dashed] (7.9,0.228) -- (8.9,0.228);
\node[font=\tiny,anchor=west] at (9.0,0.228) {$L_V{=}3w$};
\node[font=\scriptsize] at (4.5,-0.038) {(b)};
\end{scope}
\end{tikzpicture}
\caption{지연 길이 $L_{V,j}$ 가 peak 모양을 바꾸는 방식 --- 좌표는 식~\eqref{eq:lag}$\cdot$%
\eqref{eq:peakshape} 의 인과 기억 적분을 점별로 수치 평가한 값이다(방전, $w{=}1$ 단위).
(a) 단일 전이: $L_V/w=0\to0.75\to1.5\to3$ 으로 커질수록 정점이 $0.250\to0.227\to0.196\to0.149$
로 낮아지고 위치가 $0\to+0.62w\to+1.01w\to+1.50w$ 로 밀리며 꼬리가 길어진다(면적은 규격화 커널
덕에 모든 $L_V$ 에서 보존 --- 용량 $Q_j$ 불변). $L_{V}\propto|I|$ 이고 장벽 지수 탓에 저온일수록
커지므로(식~\eqref{eq:LV}$\cdot$\eqref{eq:Lqfull}), 이 한 가족이 서론 관측의 ``저온$\cdot$고율 꼬리''다.
(b) 간격 $6w$ 의 두 전이(같은 $Q$): $L_V{=}0.3w$ 에서는 두 peak 가 분리되지만(골$/$정점 $0.38$),
$L_V{=}3w$ 에서는 앞 전이의 꼬리가 이웃 peak 위에 얹혀 골이 메워지고(골$/$정점 $0.83$) 같은 $Q$
인데도 뒤 peak 가 더 높아 보인다 --- 관측의 ``꼬리가 다음 peak 와 겹친다''의 모델 안 실현이다.}
\label{fig:cand-ff2-3}
\end{figure}
